\documentclass[11pt]{article}
\usepackage{amsmath,amssymb,booktabs,geometry}
\geometry{margin=1in}

\title{Spherical FLD Relaxation Model Used by \texttt{solve\_spherical\_fld\_relaxation.py}}
\author{Quokka notes}
\date{\today}

\begin{document}
\maketitle

\section{Purpose and physical motivation}
This script computes a steady, spherically symmetric H\,II-region reference solution for a point-like ionizing source. The model is designed to match the FLD-style photon transport and local photoionization--recombination balance used in the code path being tested, while avoiding time-accurate transients.

The physical goal is to obtain the stationary ionization structure from the competition between
\begin{itemize}
\item emission of ionizing photons from a central source,
\item transport of photons by flux-limited diffusion (FLD), and
\item local absorption/recombination balance in hydrogen.
\end{itemize}

\section{Governing equations}
Let $\phi(r)$ denote the photon number density (labeled \texttt{phi} in the script, interpreted as $n_\gamma$). The steady spherical equation solved is
\begin{equation}
\frac{1}{r^2}\frac{d}{dr}\!\left(r^2 F_r\right) + S(r) - c\,\sigma_{\mathrm{HI}}\,n_{\mathrm{HI}}(\phi)\,\phi = 0,
\label{eq:steady}
\end{equation}
with FLD closure
\begin{equation}
F_r = -D\,\frac{d\phi}{dr}, \qquad D = \frac{c\,\lambda(R)}{\kappa_F},
\end{equation}
\begin{equation}
R = \frac{\left|d\phi/dr\right|}{\kappa_F\,\phi}, \qquad
\lambda(R)=\frac{2+R}{6+3R+R^2}.
\end{equation}

The transport opacity is taken as a floor value
\begin{equation}
\kappa_F = \frac{1}{L_{\mathrm{box}}},
\end{equation}
so the mean free path floor is $L_{\mathrm{box}}$.

The source is a finite-radius top-hat emitter:
\begin{equation}
S(r)=
\begin{cases}
\dfrac{3Q}{4\pi r_{\mathrm{src}}^3}, & r \le r_{\mathrm{src}},\\[4pt]
0, & r>r_{\mathrm{src}}.
\end{cases}
\end{equation}

Hydrogen chemistry is closed by local equilibrium,
\begin{equation}
c\,\sigma_{\mathrm{HI}}\,\phi\,(1-x)=\alpha_B\,n_H\,x^2,
\label{eq:xeq}
\end{equation}
where $x$ is ionized fraction and $n_{\mathrm{HI}}=n_H(1-x)$. Solving Eq.~\eqref{eq:xeq} gives
\begin{equation}
x(\phi)=\frac{2}{1+\sqrt{1+4a/b}},
\quad a=\alpha_B n_H,\quad b=c\,\sigma_{\mathrm{HI}}\,\phi,
\end{equation}
clipped to $[0,1]$.

\section{Boundary conditions}
The radial domain is $r\in[0,r_{\max}]$ with
\begin{align}
F_r(0) &= 0 \quad \text{(spherical symmetry)},\\
\phi(r_{\max}) &= 0 \quad \text{(outer Dirichlet boundary)}.
\end{align}

\section{Numerical method}
The script uses a finite-volume discretization on a 1D radial mesh.
Each nonlinear iteration does:
\begin{enumerate}
\item Picard lagging of $D$ and reaction coefficient $k_{\mathrm{reac}}=c\,\sigma_{\mathrm{HI}}\,n_{\mathrm{HI}}$ from the previous iterate.
\item Assembly of a tridiagonal linear system for the updated $\phi$.
\item Thomas-algorithm solve.
\item Under-relaxation:
\[
\phi^{(m+1)} \leftarrow (1-\omega)\phi^{(m)} + \omega\,\phi_{\mathrm{lin}}.
\]
\item Residual evaluation from Eq.~\eqref{eq:steady} and convergence test.
\end{enumerate}

The reported residual is
\begin{equation}
\mathcal{R}_{\infty} = \frac{\max_i |\mathrm{RHS}_i|}{\max_i S_i},
\end{equation}
with stopping criterion $\mathcal{R}_{\infty}<\texttt{tol}$ (or \texttt{max-iters} reached).

\section{Quokka implementation (3D, production path)}
The Python script above is a 1D spherical reference model. In Quokka, the corresponding implementation used by the H\,II-region photoionization path is a 3D Cartesian, GPU-parallel solve in
\texttt{src/particles/particle\_photoionization.hpp}, called through
\texttt{QuokkaSimulation::FillNGammaFromStromgrenAtLevel()}.

\subsection*{Where it is used in the timestep}
When \texttt{particles.use\_stromgren\_photoionization = true}, Quokka:
\begin{enumerate}
\item Reads the persistent level view of \texttt{n\_gamma\_cc} at the start of \texttt{advanceHydroAtLevel()}.
\item Passes \texttt{n\_gamma\_cc} into the Strang-split cooling source update.
\item Reuses the same \texttt{n\_gamma\_cc} for both half source steps in that hydro step.
\item After all AMR levels advance, evolves the photon field over the coarse timestep in
\texttt{AMRSimulation::ellipticSolveAllLevels()}.
\end{enumerate}
The cooling operator then converts \texttt{n\_gamma} into a local photoionization heating term
\[
\Gamma_{\mathrm{pi}} = c\,\sigma_{\mathrm{HI}}\,n_{\mathrm{HI}}\,n_\gamma\,E_{\mathrm{ion}},
\]
with the same local equilibrium closure for $x$ as Eq.~\eqref{eq:xeq}.

\subsection*{Source construction in Quokka}
Unlike the 1D top-hat source, Quokka constructs a volumetric source from stellar particles:
\begin{itemize}
\item Ionizing luminosity $Q$ is read from the stellar $Q_{\mathrm{H0}}$ lookup table (\texttt{/data/QH0} in the HDF5 table).
\item Only ionizing stellar stages contribute (high-mass non-exploding and SN progenitor stages).
\item Particle luminosities are deposited to cell centers using CIC weighting, producing a grid field \texttt{source\_q} in photons s$^{-1}$ per cell.
\item The volumetric source in the PDE is then $q_{\mathrm{vol}} = \texttt{source\_q}/\Delta V$.
\end{itemize}

\subsection*{Newton--Krylov solver differences vs. the Python script}
Quokka solves the fully implicit, time-dependent equation
\[
\frac{n_\gamma^{n+1}-n_\gamma^n}{\Delta t}
+\left[k_{\mathrm{reac}}(n_\gamma^{n+1})-\nabla\!\cdot D_{\mathrm{FLD}}(n_\gamma^{n+1})\nabla\right]n_\gamma^{n+1}=q_{\mathrm{vol}}
\]
as one composite solve over the full 3D AMR hierarchy, with
\[
k_{\mathrm{reac}} = c\,\sigma_{\mathrm{HI}}\,n_{\mathrm{HI}},
\qquad
D_{\mathrm{FLD}} = \frac{c\,\lambda(R)}{\kappa_F},
\qquad
\kappa_F = \frac{1}{L_{\mathrm{box}}}.
\]
Key implementation details:
\begin{itemize}
\item \textbf{Initialization:} One damped fixed-point predictor applies the composite frozen-coefficient operator directly, without GMRES or MLMG iteration.
\item \textbf{FLD operator:} Face-centered fluxes in each Cartesian direction with Levermore--Pomraning limiter
$\lambda=(2+R)/(6+3R+R^2)$.
\item \textbf{Nonlinear solve:} Jacobian-free Newton--Krylov with centered finite-difference Jacobian products and a projected Armijo line search that enforces nonnegative $n_\gamma$.
\item \textbf{Linear solve:} AMReX GMRES, capped at 50 iterations, with a frozen composite $\mathrm{ABecLap}$ operator as an MLMG preconditioner.
\item \textbf{Adaptive retry:} A failed nonlinear solve restores the old photon field and bisects the photon timestep recursively; successful full-timestep solves are not subcycled.
\item \textbf{Convergence metric:} the volume-weighted composite residual of the backward-Euler equation.
\end{itemize}
Thus the 1D script retains its steady tridiagonal Picard reference solve, while Quokka uses a GPU-parallel, domain-decomposed Newton--Krylov evolution of the corresponding 3D Cartesian AMR system.

\section{Parameters and script inputs}
\begin{center}
\begin{tabular}{lll}
\toprule
CLI flag & Symbol & Meaning \\
\midrule
\texttt{--Q} & $Q$ & Photon emission rate [s$^{-1}$] \\
\texttt{--alphaB} & $\alpha_B$ & Case-B recombination coefficient [cm$^3$ s$^{-1}$] \\
\texttt{--nH} & $n_H$ & Total hydrogen number density [cm$^{-3}$] \\
\texttt{--sigma-HI} & $\sigma_{\mathrm{HI}}$ & Photoionization cross section [cm$^2$] \\
\texttt{--c-light} & $c$ & Speed of light [cm s$^{-1}$] \\
\texttt{--rmax} & $r_{\max}$ & Outer radius of radial domain [cm] \\
\texttt{--r-src} & $r_{\mathrm{src}}$ & Source radius for top-hat emissivity [cm] \\
\texttt{--box-size} & $L_{\mathrm{box}}$ & Size used for $\kappa_F=1/L_{\mathrm{box}}$ [cm] \\
\texttt{--nr} & $N_r$ & Number of radial control volumes \\
\texttt{--max-iters} & --- & Maximum nonlinear iterations \\
\texttt{--tol} & --- & Residual tolerance \\
\bottomrule
\end{tabular}
\end{center}

\section{Reference Str\"omgren scale}
The script also reports the classical algebraic Str\"omgren radius,
\begin{equation}
R_S=\left(\frac{3Q}{4\pi\alpha_B n_H^2}\right)^{1/3},
\end{equation}
and a numerical $r_{50}$ where $x(r_{50})=0.5$, used for comparison diagnostics.

\end{document}
